% This file is a chapter fragment, in the same style as the other files under
% version/detailed/.  It is included by version/detailed.tex and must not
% contain a preamble of its own.
%
% Contents: the chapter in which the theorem of the cube is to be proved.  At
% present it contains only the two preparatory theorems -- the rigidity
% theorem and the seesaw theorem -- together with the cohomological inputs
% they rest on and a count of what Mathlib is missing for them.  The theorem
% of the cube itself, and its abelian-variety corollary, are deliberately not
% here yet; they are the business of a later section of this chapter, which
% will discharge the open obligation on \Cref{thm:av-cube}.
%
% Every label in this file is prefixed cube-.

\chapter{The theorem of the cube: rigidity and the seesaw theorem}
\label{cube-chapter}

This chapter is where the theorem of the cube is to be proved, in the general
form for proper varieties rather than only in the abelian-variety special case
$[n]^*L\cong L^{\otimes n^2}$.  What it contains so far are the two
preparatory theorems that the classical proof of the cube consumes:

\begin{itemize}
  \item the \emph{rigidity theorem}, \Cref{cube-thm-rigidity}: a morphism out
  of a product $X\times Y$ with $Y$ proper which collapses one fibre
  $\{x_0\}\times Y$ to a point does not depend on the second variable at all;
  \item the \emph{seesaw theorem}, \Cref{cube-thm-seesaw}: a line bundle on
  $X\times T$ with $X$ proper whose restriction to every fibre $X\times\{t\}$
  is trivial is the pullback of a line bundle on $T$, together with its
  triviality criterion \Cref{cube-cor-seesaw-trivial}.
\end{itemize}

The theorem of the cube itself is \emph{not} in this chapter yet.  Its
statement, in the form needed for abelian varieties, is \Cref{thm:av-cube} of
\Cref{chap:av-isogeny}, where it carries an open obligation; the section of
this chapter that proves it, and the derivation of
$[n]^*L\cong L^{\otimes n^2}$ from it, are to be written next and will
discharge that obligation.  Nothing in the present chapter is conditional on
the cube, and nothing in it is unproved except the three cohomological inputs
isolated in \Cref{cube-sec-inputs}, each of which is stated as its own
proposition and marked as an open obligation.

The two theorems below are the classical ones and the proofs are the classical
proofs: \cite[\S4]{MumfordAbelianVarieties} and
\cite[I, \S1]{MilneAbelianVarieties} for rigidity,
\cite[\S5]{MumfordAbelianVarieties} for the seesaw theorem, with the
cohomological inputs in the form of \cite[III, \S12]{Hartshorne}.  They are
written out in full, on the assumption that the reader has to formalise from
them and should not have to consult anything else.

\begin{remark}[Two routes to the cube]\label{cube-rem-two-routes}
The route prepared here is the classical one: rigidity, then the seesaw
theorem, then the cube.  It is not the only one.  The Stacks Project proves a
more general cube --- over an arbitrary base scheme $S$, for $X\to S$ and
$Y\to S$ flat and proper of finite presentation with geometrically integral
fibres --- as \stacksproject{0BF4}, by way of the representability of the
diagonal of the Picard functor; that proof needs deformation theory and the
theory of the Picard functor rather than semicontinuity.  The classical route
is the one chosen here because its ingredients --- global sections of a proper
variety, semicontinuity, and cohomology and base change --- are needed anyway
for the intersection-number layer on which the degree formula
\Cref{thm:av-degree-formula} rests, whereas the Picard functor is not.
\end{remark}

\section{Conventions}\label{cube-sec-conventions}

Throughout this chapter $k$ is an algebraically closed field.  A
\emph{variety} over $k$ is a $k$-scheme which is separated, of finite type and
geometrically integral over $k$; in particular it is integral, hence reduced,
irreducible, connected and non-empty.  A \emph{proper variety} is a variety
whose structure morphism to $\Spec k$ is proper.  Since $k$ is algebraically
closed and a variety is of finite type over $k$, the closed points of a
variety are exactly its $k$-rational points and they are dense in it; a
variety therefore has a $k$-rational point.  We write $X(k)$ for the set of
$k$-rational points of $X$.

For $k$-schemes $X$ and $Y$ we write $X\times Y:=X\times_kY$ and
$p_X\colon X\times Y\to X$, $p_Y\colon X\times Y\to Y$ for the two
projections.  For a point $x\in X$ we write $\{x\}\times Y:=p_X^{-1}(x)$ for
the fibre of $p_X$ over $x$; when $x$ is a $k$-rational point this fibre is
canonically isomorphic to $Y$ as a $k$-scheme, being the base change of
$Y\to\Spec k$ along $x\colon\Spec k\to X$.  For a $k$-rational point
$y_0\in Y(k)$ we write $\sigma_{y_0}\colon X\to X\times Y$ for the section of
$p_X$ determined by $\mathrm{id}_X$ and $y_0$, so that
$p_X\circ\sigma_{y_0}=\mathrm{id}_X$.

A \emph{line bundle} on a scheme $S$ is an invertible sheaf of
$\O_S$-modules, that is a sheaf of $\O_S$-modules which is locally free of
rank one.  We write $L\otimes N$ for the tensor product of two line bundles,
$L^{-1}$ for the inverse of $L$, $u^*L$ for the pullback of $L$ along a
morphism $u$, and $\mathrm{Pic}(S)$ for the group of isomorphism classes of
line bundles on $S$.  A line bundle is \emph{trivial} if it is isomorphic to
$\O_S$.  For a sheaf of $\O_S$-modules $F$ and a point $s\in S$ we write
$F(s):=F_s\otimes_{\O_{S,s}}\kappa(s)$ for the \emph{fibre} of $F$ at $s$, a
vector space over the residue field $\kappa(s)$, and $F_t$ for the
restriction of $F$ to the fibre of a morphism over a point $t$ of its target.

\begin{warning}[Composition order]\label{cube-warn-composition-order}
Composition is written in the classical order: $g\circ f$ means ``first $f$,
then $g$''.  This is the convention of \Cref{chap:av-isogeny} and it is the
opposite of the one used in the Lean sources, which write the same morphism as
\texttt{f} $\gg$ \texttt{g}.  Every composite in this chapter therefore
appears reversed in Lean.
\end{warning}

\begin{warning}[What the parallel layer supplies]
\label{cube-warn-parallel-layer}
The tensor product of sheaves of modules over a scheme, invertible sheaves and
the Picard group $\mathrm{Pic}(X)$ are being built in this project in parallel
with this chapter, so they are treated here as available and are \emph{not}
counted as missing in \Cref{cube-sec-cost}; that section counts only what sits
above them.
\end{warning}

\section{The cohomological inputs}\label{cube-sec-inputs}

This section isolates, as three separate propositions, every cohomological
fact that the two theorems of this chapter rest on.  None of the three can
currently be stated in Lean, and each is therefore marked as an open
obligation naming the theory it needs.  Nothing below hides an input inside
prose: no later proof in this chapter appeals to a cohomological fact that is
not one of these three or a consequence of them proved here.

Only the degree-zero part of the theory is needed.  The two theorems of this
chapter never mention $H^i$ for $i\geq1$: the first input is the
finite-dimensionality of the space of \emph{global sections}, and the second
and third concern $h^0$ in a family.  Higher cohomology enters only in the
\emph{proofs} of the three inputs, which are not given here, and in the Euler
characteristics that the intersection-number layer of a later chapter will
need.

\begin{proposition}[Finiteness of global sections on a proper scheme]
\label{cube-prop-coherent-finiteness}
Let $K$ be a field, let $X$ be a proper $K$-scheme and let $F$ be a coherent
sheaf of $\O_X$-modules.  Then $\Gamma(X,F)$ is a finite-dimensional
$K$-vector space.
\openobligation{AlgebraicGeometry: coherence of a sheaf of modules on a
locally noetherian scheme, and finite dimensionality of the global sections of
a coherent sheaf on a proper scheme over a field. This is the degree-zero case
of the coherence of the higher direct images of a coherent sheaf along a
proper morphism, Stacks Project Tag 02O5.}
\end{proposition}
\begin{proof}
Not proved here.  This is the case $i=0$ of \stacksproject{02O5} applied to
the structure morphism $X\to\Spec K$, whose direct images are the cohomology
$K$-vector spaces of $F$ and for which the coherent sheaves on the target are
the finite-dimensional $K$-vector spaces; see also
\cite[III, \S8]{Hartshorne}.  The degree-zero case is not appreciably easier
than the general one: the standard proof is a devissage over a finite affine
cover which computes all degrees at once.
\end{proof}

\begin{proposition}[Semicontinuity of $h^0$ in a proper flat family]
\label{cube-prop-semicontinuity}
Let $T$ be a noetherian scheme, let $f\colon W\to T$ be a proper morphism and
let $F$ be a coherent sheaf on $W$ which is flat over $T$.  For $t\in T$ let
$W_t:=W\times_T\Spec\kappa(t)$ be the fibre, let $F_t$ be the restriction of
$F$ to it, and put
\[
  h^0(t):=\dim_{\kappa(t)}\Gamma(W_t,F_t).
\]
Then $h^0$ is a well-defined function $T\to\N$ and it is upper
semicontinuous: for every $n\in\N$ the set $\{t\in T:h^0(t)\geq n\}$ is closed
in $T$.
\openobligation{AlgebraicGeometry: cohomology of a sheaf of modules on a
scheme in all degrees, flatness of a sheaf of modules over the base of a
morphism, and the semicontinuity theorem for the fibrewise cohomology of a
coherent sheaf in a proper flat family.}
\end{proposition}
\begin{proof}
\uses{cube-prop-coherent-finiteness}
Not proved here, apart from the finiteness of the values.  That $h^0(t)$ is a
natural number is \Cref{cube-prop-coherent-finiteness} applied to the proper
$\kappa(t)$-scheme $W_t$ and the coherent sheaf $F_t$.  The semicontinuity
statement is \cite[III, Theorem 12.8]{Hartshorne}; its proof computes the
cohomology of all the fibres from a single complex of flat modules over the
base, so it needs $H^i$ in every degree and not only $H^0$.
\end{proof}

\begin{proposition}[Cohomology and base change in degree zero, after Grauert]
\label{cube-prop-grauert}
Let $T$ be an integral noetherian scheme, let $f\colon W\to T$ be a proper
flat morphism and let $F$ be a coherent sheaf on $W$ which is flat over $T$.
Suppose that the function $h^0$ of \Cref{cube-prop-semicontinuity} is constant
equal to $r$ on $T$.  Then:
\begin{enumerate}
  \item $f_*F$ is a locally free sheaf of $\O_T$-modules of rank $r$;
  \item for every $t\in T$ the canonical base change map
  \[
    (f_*F)(t)\longrightarrow\Gamma(W_t,F_t)
  \]
  from the fibre of $f_*F$ at $t$ to the global sections of $F_t$ is an
  isomorphism of $\kappa(t)$-vector spaces;
  \item the restriction to $W_t$ of the counit
  $\varepsilon\colon f^*f_*F\to F$ of the adjunction $f^*\dashv f_*$ is, under
  the identification of (2), the evaluation map
  \[
    \Gamma(W_t,F_t)\otimes_{\kappa(t)}\O_{W_t}\longrightarrow F_t
  \]
  sending $s\otimes1$ to $s$.
\end{enumerate}
\openobligation{AlgebraicGeometry: cohomology of a sheaf of modules on a
scheme, the base change map for the pushforward of a coherent sheaf along a
proper morphism, and cohomology and base change in the form of the theorem of
Grauert for a proper flat morphism over an integral base.}
\end{proposition}
\begin{proof}
\uses{cube-prop-semicontinuity, cube-prop-coherent-finiteness}
Not proved here.  Parts (1) and (2) are Grauert's theorem,
\cite[III, Theorem 12.11]{Hartshorne}, whose hypotheses are exactly those
above; see also \cite[\S5]{MumfordAbelianVarieties}, where the statement is
given in the form used below.  Part (3) is a compatibility rather than an
extra input: the base change map of (2) is by construction the map induced on
fibres by $\varepsilon$, so once (2) is known, (3) says that the counit of an
adjunction is compatible with base change.  It is recorded here as part of the
statement so that the proof of \Cref{cube-thm-seesaw} does not have to appeal
to it silently.
\end{proof}

The two consequences of \Cref{cube-prop-coherent-finiteness} that the rest of
the chapter uses are the following.  Both are stated over the algebraically
closed field $k$, which is all that is needed; \Cref{cube-rem-general-field}
says what changes over a general field.

\begin{lemma}[Global sections of a proper variety]
\label{cube-lem-global-sections}
Let $X$ be a proper integral scheme over the algebraically closed field $k$
--- in particular, a proper variety over $k$.  Then the structure map
$k\to\Gamma(X,\O_X)$ is an isomorphism; that is, $\Gamma(X,\O_X)=k$ and
$\dim_k\Gamma(X,\O_X)=1$.
\end{lemma}
\begin{proof}
\uses{cube-prop-coherent-finiteness}
Put $A:=\Gamma(X,\O_X)$.  The structure sheaf is coherent, so
\Cref{cube-prop-coherent-finiteness} makes $A$ a finite-dimensional
$k$-algebra.  Since $X$ is integral, $A$ is a subring of the function field of
$X$ and is therefore an integral domain.  A finite-dimensional algebra over a
field which is an integral domain is a field: multiplication by a non-zero
element is an injective $k$-linear endomorphism of a finite-dimensional
$k$-vector space, hence surjective, so the element is invertible.  Thus $A$ is
a finite field extension of $k$, and $k$ is algebraically closed, so $A=k$.
\end{proof}

\begin{remark}[The same statement over a general field]
\label{cube-rem-general-field}
Over a field $K$ that is not algebraically closed the conclusion
$\Gamma(X,\O_X)=K$ needs $X$ to be \emph{geometrically} integral and not
merely integral --- for $X=\Spec L$ with $L/K$ a non-trivial finite extension,
$X$ is proper and integral and $\Gamma(X,\O_X)=L$.  With geometric
integrality the statement is true, and is \stacksproject{0BUG}, part (8); see
also \stacksproject{0FD2}.  Its proof needs one input beyond
\Cref{cube-prop-coherent-finiteness}, namely that the formation of global
sections commutes with the flat base change $K\to\overline K$, which for a
quasi-compact separated scheme is elementary --- tensoring the equaliser
defining $\Gamma$ over a finite affine cover with the flat $K$-module
$\overline K$ --- but is still a statement Mathlib does not have.  Nothing in
this chapter needs it, because every scheme to which
\Cref{cube-lem-global-sections} is applied below is a proper variety over
$k$ itself or a fibre over a $k$-rational point, and such a fibre again has
$k$ as its residue field.  This is why \Cref{cube-lem-h0-constant} is arranged
to compute $h^0$ only at closed points and to reach the remaining points of
the base by semicontinuity.
\end{remark}

\begin{lemma}[Sections of a trivial line bundle on a proper variety]
\label{cube-lem-nowhere-vanishing}
Let $X$ be a proper integral scheme over $k$ and let $L$ be a trivial line
bundle on $X$.  Then $\dim_k\Gamma(X,L)=1$, and every non-zero global section
$s\in\Gamma(X,L)$ generates $L$ at every point of $X$; equivalently, the image
of $s$ in the fibre $L(x)$ is non-zero for every $x\in X$.
\end{lemma}
\begin{proof}
\uses{cube-lem-global-sections}
Both assertions are invariant under isomorphism of line bundles, so we may
choose an isomorphism $L\cong\O_X$ and assume $L=\O_X$.  Then
$\Gamma(X,L)=\Gamma(X,\O_X)=k$ by \Cref{cube-lem-global-sections}, which gives
the dimension.  A non-zero section is then a non-zero element $c\in k$, hence
a unit of $\Gamma(X,\O_X)$ and so a unit in every local ring $\O_{X,x}$; a
unit generates $\O_{X,x}$ as a module over itself, and its image in
$\O_{X,x}/\m_x=\kappa(x)$ is non-zero.
\end{proof}

\section{Rigidity}\label{cube-sec-rigidity}

The rigidity theorem says that a morphism out of a product cannot collapse one
fibre of a projection whose fibre direction is proper without collapsing all
of them.  The two steps of the proof are recorded separately: the properness
step, which produces an open subset of the base on which the collapse already
happens, and the density step, which propagates it to the whole source.

\begin{lemma}[A morphism from a proper variety to an affine scheme is
constant]\label{cube-lem-constant-to-affine}
Let $Y$ be a proper integral scheme over $k$ and let $g\colon Y\to V$ be a
morphism of $k$-schemes with $V$ affine.  Then $g$ factors as
\[
  Y\xrightarrow{\ \pi\ }\Spec k\xrightarrow{\ v\ }V
\]
for a unique $k$-rational point $v\in V(k)$, where $\pi$ is the structure
morphism.  In particular the set-theoretic image of $g$ is the single point
$v$.
\end{lemma}
\begin{proof}
\uses{cube-lem-global-sections}
Write $V=\Spec A$ with $A$ a $k$-algebra.  By the adjunction between the
global sections functor and $\Spec$, morphisms of $k$-schemes $Y\to\Spec A$
correspond bijectively to morphisms of $k$-algebras $A\to\Gamma(Y,\O_Y)$.  By
\Cref{cube-lem-global-sections} we have $\Gamma(Y,\O_Y)=k$, so such a
morphism is a $k$-algebra map $A\to k$, that is a $k$-rational point $v$ of
$V$; and the morphism $Y\to\Spec A$ it corresponds to is by construction the
composite of $\pi$ with $v$.  Uniqueness of $v$ is the injectivity half of the
same bijection.
\end{proof}

\begin{lemma}[Properness makes the collapsing locus open]
\label{cube-lem-collapse-neighbourhood}
Let $X$ and $Z$ be $k$-schemes, let $Y$ be a proper $k$-scheme and let
$f\colon X\times Y\to Z$ be a morphism.  Let $x_0\in X$ and let $U\subseteq Z$
be an open subset with $f(\{x_0\}\times Y)\subseteq U$.  Then there is an open
subset $V\subseteq X$ with $x_0\in V$ and
\[
  f\bigl(p_X^{-1}(V)\bigr)\subseteq U .
\]
\end{lemma}
\begin{proof}
The projection $p_X\colon X\times Y\to X$ is the base change of
$Y\to\Spec k$ along $X\to\Spec k$, so it is proper, hence universally closed,
hence a closed map of the underlying topological spaces.

Put $W:=f^{-1}(Z\setminus U)$, a closed subset of $X\times Y$.  Then
$p_X(W)$ is closed in $X$, and $x_0\notin p_X(W)$: a point of
$W\cap p_X^{-1}(x_0)$ would be a point of $\{x_0\}\times Y$ whose image under
$f$ lies outside $U$, contradicting the hypothesis.  So
$V:=X\setminus p_X(W)$ is an open subset of $X$ containing $x_0$, and
$p_X^{-1}(V)$ is disjoint from $W$, which is exactly the assertion
$f(p_X^{-1}(V))\subseteq U$.
\end{proof}

\begin{theorem}[Rigidity]\label{cube-thm-rigidity}
Let $X$ be a variety over $k$, let $Y$ be a proper variety over $k$ and let
$Z$ be a separated $k$-scheme, locally of finite type over $k$.  Let
$f\colon X\times Y\to Z$ be a morphism of $k$-schemes and suppose that for
some point $x_0\in X$ the image $f(\{x_0\}\times Y)$ is a single point of $Z$.
Then there is a unique morphism $h\colon X\to Z$ with
\[
  f=h\circ p_X ,
\]
namely $h=f\circ\sigma_{y_0}$ for any $y_0\in Y(k)$.  In particular $f$ does
not depend on the second variable.
\end{theorem}
\begin{proof}
\uses{cube-lem-constant-to-affine, cube-lem-collapse-neighbourhood}
Uniqueness is immediate.  The variety $Y$ has a $k$-rational point $y_0$, and
$\sigma_{y_0}$ is a section of $p_X$, so $f=h\circ p_X$ forces
$h=h\circ p_X\circ\sigma_{y_0}=f\circ\sigma_{y_0}$.  Fix such a $y_0$ and put
\[
  h:=f\circ\sigma_{y_0}\colon X\to Z .
\]
It remains to prove $f=h\circ p_X$.

\emph{Step 1: an open set on which the collapse happens.}  Let $z_0$ be the
point with $f(\{x_0\}\times Y)=\{z_0\}$ and choose an affine open subset
$U\subseteq Z$ containing $z_0$.  By
\Cref{cube-lem-collapse-neighbourhood} there is an open subset $V\subseteq X$
with $x_0\in V$ and $f(p_X^{-1}(V))\subseteq U$.  A morphism whose
set-theoretic image lies in an open subset factors through that open subset as
an open subscheme, so the restriction of $f$ to $p_X^{-1}(V)$ factors through
$U$.

\emph{Step 2: $f$ and $h\circ p_X$ agree at the $k$-rational points of
$p_X^{-1}(V)$.}  Let $x\in V(k)$ and $y\in Y(k)$.  The fibre
$\{x\}\times Y$ is isomorphic to $Y$ as a $k$-scheme, because $x$ is a
$k$-rational point, and by Step 1 the restriction of $f$ to it factors through
the affine open $U$.  So \Cref{cube-lem-constant-to-affine}, applied to the
proper integral $k$-scheme $Y$ and the morphism $\{x\}\times Y\to U$, shows
that this restriction is constant.  Its value at
$(x,y_0)=\sigma_{y_0}(x)$ is $h(x)$, so its value at $(x,y)$ is $h(x)$ as
well; that is, $f(x,y)=h(x)=(h\circ p_X)(x,y)$.  Every $k$-rational point of
$p_X^{-1}(V)$ is of this form, so $f$ and $h\circ p_X$ agree at all of them.

\emph{Step 3: propagation by density.}  The scheme $X\times Y$ is a variety
over $k$, being a product of two geometrically integral $k$-schemes of finite
type; in particular it is irreducible and reduced.  The subset
$p_X^{-1}(V)=V\times Y$ is open in $X\times Y$ and non-empty, hence dense.
Its closed points are its $k$-rational points, because $k$ is algebraically
closed and $X\times Y$ is of finite type over $k$.  By Step 2 the two
morphisms $f$ and $h\circ p_X$ agree at every closed point of the dense open
subset $p_X^{-1}(V)$; both are morphisms of $k$-schemes into the separated
$k$-scheme $Z$, and $X\times Y$ is reduced.  Hence $f=h\circ p_X$.
\end{proof}

\begin{remark}[The Lean shape of the proof]\label{cube-rem-rigidity-mathlib}
Steps 1 and 3 are already available in Mathlib \texttt{v4.33.1}.  For Step 3,
the statement that two $K$-morphisms out of a reduced $K$-scheme locally of
finite type into a separated $K$-scheme locally of finite type are equal as
soon as they agree at the closed points of a dense locally closed subset, for
$K$ algebraically closed, is \texttt{AlgebraicGeometry.ext\_of\_apply\_eq} in
\texttt{Mathlib/AlgebraicGeometry/AlgClosed/Basic.lean}; the identification of
closed points with $k$-rational points is
\texttt{AlgebraicGeometry.pointEquivClosedPoint} in the same file, and their
density is the Jacobson property of a scheme locally of finite type over a
field.  For Step 1, properness is stable under base change in Mathlib and
\texttt{AlgebraicGeometry.Scheme.Hom.isClosedMap} turns a universally closed
morphism into a closed map.  The only ingredient of the whole proof that is
not available is \Cref{cube-lem-constant-to-affine}, and only because of
\Cref{cube-lem-global-sections} inside it.
\end{remark}

\begin{warning}[Properness is needed on the factor that gets collapsed]
\label{cube-warn-proper-factor}
In \Cref{cube-thm-rigidity} the properness hypothesis is on $Y$, the factor
whose fibre $\{x_0\}\times Y$ is assumed to collapse, and not on $X$.  It
cannot be moved: with $X$ proper and $Y$ merely connected the statement is
false.  For a counterexample take $X=Z=\mathbf P^1_k$ and
$Y=\mathbf A^1_k$ and let
\[
  f\colon\mathbf P^1_k\times\mathbf A^1_k\to\mathbf P^1_k,\qquad
  f\bigl([u:v],t\bigr)=[u:v+tu] ,
\]
which is a morphism because $u=0$ forces $v\neq0$ and then $v+tu=v\neq0$.
Take $x_0=[0:1]$.  Then $f(\{x_0\}\times Y)=\{[0:1]\}$ is a single point, $X$
is proper and connected and $Y$ is connected; but $f([1:0],t)=[1:t]$ depends
on $t$, so $f$ does not factor through $p_X$.  What properness of $Y$ is used
for is exactly that $p_X$ is a closed map, in
\Cref{cube-lem-collapse-neighbourhood}.  In the application to the theorem of
the cube both factors are proper varieties, so the distinction is invisible
there; it matters when the statement is formalised, because the hypothesis
would otherwise be attached to the wrong variable.
\end{warning}

\begin{remark}[Weaker hypotheses]\label{cube-rem-rigidity-weaker}
Irreducibility of $X$ is used only for the density of $V\times Y$ in Step 3,
and reducedness of $X\times Y$ only in the last sentence of that step; with
$X$ merely connected and reduced the conclusion still holds, propagating the
collapse along a connected chain of irreducible components, but the argument
is longer and the version above is all that the theorem of the cube needs.
Integrality of $Y$, by contrast, is essential: it is what makes
$\Gamma(Y,\O_Y)=k$ in \Cref{cube-lem-global-sections}, and hence what makes a
morphism from $Y$ to an affine scheme constant.
\end{remark}

\section{The seesaw theorem}\label{cube-sec-seesaw}

The seesaw theorem is stated here in the form the theorem of the cube will
consume: triviality on every fibre of a projection whose fibre direction is
proper implies that the bundle comes from the base.  Three auxiliary
statements come first: that the projection is a proper flat family, a
fibrewise criterion for a morphism of line bundles to be an isomorphism, and
the constancy of $h^0$.

\begin{lemma}[The projection is a proper flat family]
\label{cube-lem-projection-proper-flat}
Let $X$ be a proper variety over $k$ and let $T$ be a variety over $k$.  Then:
\begin{enumerate}
  \item $X\times T$ is a variety over $k$; in particular it is an integral
  noetherian scheme;
  \item $p_T\colon X\times T\to T$ is proper, flat and of finite type, and it
  is a closed map;
  \item for $t\in T$ the fibre of $p_T$ over $t$, written $X_t$, is the base
  change $X\times_k\Spec\kappa(t)$, a proper $\kappa(t)$-scheme; for
  $t\in T(k)$ it is isomorphic to $X$ as a $k$-scheme;
  \item every line bundle $L$ on $X\times T$ is coherent and flat over $T$.
\end{enumerate}
\end{lemma}
\begin{proof}
For (1): a product of two geometrically integral $k$-schemes is geometrically
integral, and separatedness and finite type are stable under products; a
scheme of finite type over a field is noetherian.  For (2): $p_T$ is the base
change of $X\to\Spec k$ along $T\to\Spec k$, and properness, flatness and
finite type are all stable under base change --- flatness because every
morphism to the spectrum of a field is flat.  A proper morphism is
universally closed and hence a closed map.  For (3): the fibre of a base
change is the base change of the fibre, so the fibre of $p_T$ over $t$ is
$X\times_k\Spec\kappa(t)$, which is proper over $\kappa(t)$ because
properness is stable under base change; for $t\in T(k)$ we have
$\kappa(t)=k$ and the fibre is $X$ itself.  For (4): $L$ is locally free of
rank one, hence coherent on the noetherian scheme $X\times T$, and a locally
free sheaf is flat over $X\times T$, hence flat over $T$ because $p_T$ is flat
by (2).
\end{proof}

\begin{lemma}[Fibrewise criterion for a morphism of line bundles]
\label{cube-lem-invertible-fibrewise}
Let $S$ be a scheme, let $M$ and $N$ be line bundles on $S$ and let
$\varphi\colon M\to N$ be a morphism of $\O_S$-modules.  Suppose that for
every point $s\in S$ the induced $\kappa(s)$-linear map on fibres
$\varphi(s)\colon M(s)\to N(s)$ is non-zero.  Then $\varphi$ is an
isomorphism.  Moreover, without any hypothesis on $\varphi$, the set
$\{s\in S:\varphi(s)=0\}$ is closed in $S$.
\end{lemma}
\begin{proof}
Being an isomorphism is local on $S$, so it suffices to treat an open subset
$U\subseteq S$ on which both $M$ and $N$ are free of rank one; let $m$ and $n$
be generators.  Then $\varphi(m)=a\cdot n$ for a unique
$a\in\Gamma(U,\O_U)$, and for $s\in U$ the map $\varphi(s)$ sends the basis
vector $\overline m$ of the one-dimensional space $M(s)$ to
$\overline a\cdot\overline n$, where $\overline a$ is the image of $a$ in
$\kappa(s)$.  So $\varphi(s)=0$ if and only if $\overline a=0$, that is
$a\in\m_s$; the vanishing locus of $\varphi$ on $U$ is therefore the zero set
of the section $a$, which is closed, and this proves the last assertion since
such opens cover $S$.

Under the hypothesis, $a\notin\m_s$ for every $s\in U$, that is $a$ is a unit
in every local ring $\O_{S,s}$ with $s\in U$.  A section of $\O_U$ which is a
unit in every local ring of $U$ is a unit of $\Gamma(U,\O_U)$: it is
invertible on an open cover, and the local inverses agree on overlaps because
inverses are unique, so they glue.  Hence $\varphi|_U$ is the isomorphism
sending $m$ to $n$ followed by multiplication by a unit, and is therefore an
isomorphism.
\end{proof}

\begin{lemma}[$h^0$ is identically one]\label{cube-lem-h0-constant}
Let $X$ be a proper variety over $k$, let $T$ be a variety over $k$ and let
$L$ be a line bundle on $X\times T$ whose restriction $L_t$ to
$X\times\{t\}$ is trivial for every $t\in T(k)$.  Then the function $h^0$ of
\Cref{cube-prop-semicontinuity}, formed for $p_T\colon X\times T\to T$ and
$F=L$, is constant equal to $1$ on all of $T$, and not only at the
$k$-rational points.
\end{lemma}
\begin{proof}
\uses{cube-prop-semicontinuity, cube-lem-nowhere-vanishing,
cube-lem-projection-proper-flat}
By \Cref{cube-lem-projection-proper-flat} the morphism $p_T$ is proper and
$L$ is coherent and flat over $T$, and $T$ is noetherian, so $h^0$ is defined
and upper semicontinuous by \Cref{cube-prop-semicontinuity}.

For $t\in T(k)$ the fibre is $X$ and $L_t$ is trivial by hypothesis, so
$h^0(t)=\dim_k\Gamma(X,L_t)=1$ by \Cref{cube-lem-nowhere-vanishing}.  Thus
$h^0$ takes the value $1$ at every closed point of $T$.

Now let $t\in T$ be arbitrary.  The set $\{h^0\geq2\}$ is closed by
semicontinuity; if it contained $t$ it would contain the closure of $\{t\}$,
and that closure contains a closed point of $T$, because $T$ is of finite type
over a field and hence a Jacobson scheme.  That contradicts $h^0=1$ at closed
points, so $h^0(t)\leq1$.  The set $\{h^0\geq1\}$ is closed as well and
contains every closed point of $T$; the closed points are dense, so this set
is all of $T$ and $h^0(t)\geq1$.  Hence $h^0(t)=1$.
\end{proof}

\begin{theorem}[Seesaw theorem]\label{cube-thm-seesaw}
Let $X$ be a proper variety over $k$, let $T$ be a variety over $k$ and let
$L$ be a line bundle on $X\times T$ whose restriction to $X\times\{t\}$ is
trivial for every $t\in T(k)$.  Then
\[
  M:=p_{T*}L
\]
is a line bundle on $T$ and the counit of the adjunction
$p_T^*\dashv p_{T*}$ is an isomorphism
\[
  \varepsilon\colon p_T^*M\;\xrightarrow{\ \sim\ }\;L .
\]
In particular $L$ is the pullback along $p_T$ of a line bundle on $T$.
\end{theorem}
\begin{proof}
\uses{cube-lem-h0-constant, cube-prop-grauert,
cube-lem-invertible-fibrewise, cube-lem-nowhere-vanishing,
cube-lem-projection-proper-flat}
By \Cref{cube-lem-projection-proper-flat} the morphism
$p_T\colon X\times T\to T$ is proper and flat, $L$ is coherent and flat over
$T$, and $T$ is an integral noetherian scheme.  By
\Cref{cube-lem-h0-constant} the function $h^0$ is constant equal to $1$ on
$T$.  So \Cref{cube-prop-grauert} applies with $r=1$ and yields three things:
$M=p_{T*}L$ is locally free of rank one, that is a line bundle; for every
$t\in T$ the base change map $M(t)\to\Gamma(X_t,L_t)$ is an isomorphism of
one-dimensional $\kappa(t)$-vector spaces; and the restriction of
$\varepsilon\colon p_T^*M\to L$ to the fibre over $t$ is, under that
identification, the evaluation map
\[
  \Gamma(X_t,L_t)\otimes_{\kappa(t)}\O_{X_t}\longrightarrow L_t .
\]

It remains to prove that $\varepsilon$, which is a morphism of line bundles on
$X\times T$, is an isomorphism.  By \Cref{cube-lem-invertible-fibrewise} it is
enough to show that $\varepsilon(w)\neq0$ for every point
$w\in X\times T$.  Let
\[
  N:=\{w\in X\times T:\varepsilon(w)=0\} ,
\]
which is a closed subset of $X\times T$ by the last assertion of
\Cref{cube-lem-invertible-fibrewise}.

First, $N$ misses every fibre over a $k$-rational point.  Indeed, let
$t\in T(k)$; then $X_t\cong X$ is a proper integral $k$-scheme and $L_t$ is
trivial by hypothesis, so by \Cref{cube-lem-nowhere-vanishing} the space
$\Gamma(X_t,L_t)$ is one-dimensional and any basis vector $s$ of it is nowhere
vanishing on $X_t$.  By the description of the restriction of $\varepsilon$
above, for $w\in X_t$ the map $\varepsilon(w)$ sends the basis vector
$s\otimes1$ of the one-dimensional space $(p_T^*M)(w)$ to the image of $s$ in
$L_t(w)$, which is non-zero.  So $\varepsilon(w)\neq0$ and
$w\notin N$.

Second, $p_T$ is a closed map by \Cref{cube-lem-projection-proper-flat}, so
$p_T(N)$ is a closed subset of $T$.  By the previous paragraph it contains no
$k$-rational point of $T$, hence no closed point of $T$.  A non-empty closed
subset of $T$ contains a closed point of $T$, because $T$ is of finite type
over a field and hence a Jacobson scheme; therefore $p_T(N)=\emptyset$ and
so $N=\emptyset$.

Thus $\varepsilon(w)\neq0$ for every $w\in X\times T$, and
\Cref{cube-lem-invertible-fibrewise} shows that $\varepsilon$ is an
isomorphism.
\end{proof}

\begin{corollary}[Seesaw, triviality form]\label{cube-cor-seesaw-trivial}
Let $X$ be a proper variety over $k$, let $T$ be a variety over $k$ and let
$L$ be a line bundle on $X\times T$.  Suppose that
\begin{enumerate}
  \item the restriction of $L$ to $X\times\{t\}$ is trivial for every
  $t\in T(k)$, and
  \item the restriction of $L$ to $\{x_0\}\times T$ is trivial for some
  $x_0\in X(k)$.
\end{enumerate}
Then $L$ is trivial.
\end{corollary}
\begin{proof}
\uses{cube-thm-seesaw}
By \Cref{cube-thm-seesaw} there are a line bundle $M$ on $T$ and an
isomorphism $p_T^*M\cong L$.  Let $\sigma\colon T\to X\times T$ be the section
of $p_T$ determined by $x_0$ and $\mathrm{id}_T$, so that
$p_T\circ\sigma=\mathrm{id}_T$ and $\sigma$ identifies $T$ with
$\{x_0\}\times T$.  Pulling the isomorphism back along $\sigma$ gives
\[
  M\cong(p_T\circ\sigma)^*M\cong\sigma^*p_T^*M\cong\sigma^*L ,
\]
and $\sigma^*L$ is the restriction of $L$ to $\{x_0\}\times T$, which is
trivial by hypothesis (2).  So $M$ is trivial and therefore
$L\cong p_T^*M\cong p_T^*\O_T=\O_{X\times T}$.
\end{proof}

\begin{remark}[Why the base has to be a variety]\label{cube-rem-seesaw-base}
Integrality of $T$ is used twice in \Cref{cube-thm-seesaw} and cannot simply
be dropped.  It is a hypothesis of \Cref{cube-prop-grauert}: over a
non-reduced base a constant $h^0$ does not force the pushforward to be locally
free.  And it is used in \Cref{cube-lem-h0-constant}, where the hypothesis at
the $k$-rational points is propagated to the remaining points of $T$ using
that the closed points are dense.  For the theorem of the cube this costs
nothing, since the parameter spaces occurring there are varieties.
\end{remark}

\begin{remark}[What the cube will use]\label{cube-rem-what-the-cube-uses}
The theorem of the cube, to be proved in the next section of this chapter,
will use \Cref{cube-thm-rigidity} to see that a suitable difference of
pullbacks of a line bundle is trivial on the relevant coordinate slices, and
\Cref{cube-cor-seesaw-trivial} --- with $T$ the third factor and $X$ the
product of the first two --- to conclude that it is trivial on the whole
triple product.  Its abelian-variety corollary, $[n]^*L\cong L^{\otimes n^2}$
for symmetric $L$, is \Cref{thm:av-mulby-pullback} of
\Cref{chap:av-isogeny}, and the open obligation on \Cref{thm:av-cube} there is
what that section will discharge.
\end{remark}

\section{Formalisation cost}\label{cube-sec-cost}

This section counts what Mathlib \texttt{v4.33.1} does not have and the two
theorems of this chapter need.  It is scoped to the chapter as it stands, that
is to \Cref{cube-thm-rigidity}, \Cref{cube-thm-seesaw} and
\Cref{cube-cor-seesaw-trivial} together with the statements they use.  The
section of this chapter that proves the theorem of the cube will extend it,
and the intersection-number layer behind \Cref{thm:av-degree-formula} is not
counted here at all.

As said in \Cref{cube-warn-parallel-layer}, the tensor product of sheaves of
modules over a scheme, invertible sheaves and the Picard group of a scheme are
being built in this project in parallel and are not counted as missing; what
follows counts only what sits above them.

\subsection*{What Mathlib \texttt{v4.33.1} already has}

So that the count below is honest, here is what is available and used.
\begin{enumerate}
  \item \texttt{AlgebraicGeometry.Scheme.Modules}, the abelian category of
  sheaves of modules over a scheme, in
  \texttt{Mathlib/AlgebraicGeometry/Modules/Sheaf.lean}, with
  \texttt{pullback} and \texttt{pushforward}, their adjunction --- whose
  counit is the map $\varepsilon$ of \Cref{cube-thm-seesaw} --- and full
  pseudofunctoriality.
  \item \texttt{SheafOfModules.IsLocallyFree} and
  \texttt{SheafOfModules.IsQuasicoherent}, in
  \texttt{Mathlib/Algebra/Category/ModuleCat/Sheaf/LocallyFree.lean} and
  \texttt{Quasicoherent.lean}; a monoidal structure on \emph{presheaves} of
  modules and their sheafification; and \texttt{CommRing.Pic} together with
  \texttt{Module.Invertible} for a ring, in
  \texttt{Mathlib/RingTheory/PicardGroup.lean}.
  \item \texttt{CategoryTheory.Sheaf.H}, the cohomology of an \emph{abelian}
  sheaf on an arbitrary site, defined as $\mathrm{Ext}$ from the constant
  sheaf, with the Cech and Mayer--Vietoris machinery, in
  \texttt{Mathlib/CategoryTheory/Sites/SheafCohomology/}; and the Euler
  characteristic of a homological complex in
  \texttt{Mathlib/Algebra/Homology/EulerCharacteristic.lean}.  Neither is
  applied to schemes or to sheaves of modules anywhere in Mathlib, which is
  precisely the gap the three open obligations sit in.
  \item Everything topological and morphism-theoretic that the proofs above
  use: \texttt{IsProper} and its stability under base change,
  \texttt{Scheme.Hom.isClosedMap} for a universally closed morphism,
  \texttt{Flat} and its stability under base change,
  \texttt{GeometricallyIntegral} with the instances for the two projections of
  a product, the Jacobson property of a scheme locally of finite type over a
  field together with the density of closed points,
  \texttt{AlgebraicGeometry.pointEquivClosedPoint} and
  \texttt{AlgebraicGeometry.ext\_of\_apply\_eq} over an algebraically closed
  base field, \texttt{AlgebraicGeometry.ext\_of\_isDominant}, the adjunction
  between global sections and $\Spec$, and \texttt{Scheme.residueField}.
\end{enumerate}

\subsection*{Definitions missing: six, of which four carry content}

\begin{enumerate}
  \item \emph{Coherence of a sheaf of modules on a locally noetherian
  scheme}, \texttt{SheafOfModules.IsCoherent}, with the facts that it is
  stable under pullback and that a locally free sheaf of finite rank is
  coherent.  Needed to state \Cref{cube-prop-coherent-finiteness},
  \Cref{cube-prop-semicontinuity} and \Cref{cube-prop-grauert}, and used in
  \Cref{cube-lem-projection-proper-flat}.  Would live in
  \texttt{WeilConjectures/Mathlib/Algebra/Category/ModuleCat/Sheaf/Coherent.lean},
  mirroring
  \texttt{Mathlib/Algebra/Category/ModuleCat/Sheaf/Quasicoherent.lean}.
  \item \emph{The fibre of a sheaf of modules at a point}, $F(s)$, as a vector
  space over the residue field, with its functoriality in $F$.  Needed to
  state \Cref{cube-lem-invertible-fibrewise} and part (2) of
  \Cref{cube-prop-grauert}, and used in
  \Cref{cube-lem-nowhere-vanishing} and \Cref{cube-thm-seesaw}.  Would live in
  \texttt{WeilConjectures/Mathlib/AlgebraicGeometry/Modules/Sheaf.lean}.
  \item \emph{The base change map} $(f_*F)(t)\to\Gamma(W_t,F_t)$.  Needed to
  state part (2) of \Cref{cube-prop-grauert} and used in
  \Cref{cube-thm-seesaw}.  Would live in a new file
  \texttt{WeilConjectures/AlgebraicGeometry/Cohomology/BaseChange.lean}, which
  mirrors no Mathlib file.
  \item \emph{The fibrewise dimension function} $h^0\colon T\to\N$,
  $t\mapsto\dim_{\kappa(t)}\Gamma(W_t,F_t)$, which needs
  \Cref{cube-prop-coherent-finiteness} to be well behaved.  Needed to state
  \Cref{cube-prop-semicontinuity}, \Cref{cube-prop-grauert} and
  \Cref{cube-lem-h0-constant}.  Same file as item (3).
  \item \emph{The restriction of a sheaf of modules to a fibre of a
  morphism}, $F_t$ on $W_t$.  This is glue rather than content: it is the
  pullback functor of
  \texttt{Mathlib/AlgebraicGeometry/Modules/Sheaf.lean} along the fibre
  inclusion of \texttt{Mathlib/AlgebraicGeometry/Fiber.lean}, so it is an
  abbreviation.  Needed by every statement of \Cref{cube-sec-inputs} and
  \Cref{cube-sec-seesaw}.
  \item \emph{Variety over a field}: an abbreviation for separated, of finite
  type and geometrically integral, with the derived instances of integrality,
  irreducibility and connectedness.  Also glue --- all three predicates are in
  Mathlib.  Needed by every statement of this chapter.  Would live in
  \texttt{WeilConjectures/Mathlib/AlgebraicGeometry/Geometrically/Integral.lean}.
\end{enumerate}

\subsection*{Statements missing: three left open, ten proved here}

The three inputs left open are \Cref{cube-prop-coherent-finiteness},
\Cref{cube-prop-semicontinuity} and \Cref{cube-prop-grauert}, all in
\Cref{cube-sec-inputs}; they are the only unproved statements in the chapter.
Beyond the definitions above, each needs a theory Mathlib does not have.
\begin{enumerate}
  \item \Cref{cube-prop-coherent-finiteness} needs the coherence of the direct
  images of a coherent sheaf along a proper morphism, \stacksproject{02O5}, of
  which it is the degree-zero case; the standard proof needs the cohomology of
  a sheaf of modules on a scheme in all degrees and a devissage over a finite
  affine cover.  It is used by \Cref{cube-lem-global-sections} and hence,
  transitively, by both theorems of the chapter, and it is the \emph{only}
  missing input of \Cref{cube-thm-rigidity}.  A new file
  \texttt{WeilConjectures/AlgebraicGeometry/Cohomology/Proper.lean}.
  \item \Cref{cube-prop-semicontinuity} needs, in addition, flatness of a
  sheaf of modules over the base of a morphism and the semicontinuity theorem,
  \cite[III, Theorem 12.8]{Hartshorne}.  Used by
  \Cref{cube-lem-h0-constant}.  A new file
  \texttt{WeilConjectures/AlgebraicGeometry/Cohomology/Semicontinuity.lean}.
  \item \Cref{cube-prop-grauert} needs cohomology and base change in all
  degrees, \cite[III, Theorem 12.11]{Hartshorne}.  Used by
  \Cref{cube-thm-seesaw}.  Would live in the \texttt{BaseChange.lean} file of
  item (3) of the previous list.
\end{enumerate}
There is one further definition that no statement of this chapter mentions but
that the proofs of those three obligations would need: the cohomology
$H^i(X,F)$ of a sheaf of modules on a scheme as a module over
$\Gamma(X,\O_X)$, with its Cech computation for a quasi-coherent sheaf on a
separated scheme and the base change spectral sequence.  It is the same
definition that the Euler characteristic $\chi(F)$ of the intersection-number
layer needs, so it is a shared cost and is not charged to this chapter.

Ten statements are proved here in full and are new to Lean:
\Cref{cube-lem-global-sections},
\Cref{cube-lem-nowhere-vanishing},
\Cref{cube-lem-constant-to-affine},
\Cref{cube-lem-collapse-neighbourhood},
\Cref{cube-thm-rigidity},
\Cref{cube-lem-projection-proper-flat},
\Cref{cube-lem-invertible-fibrewise},
\Cref{cube-lem-h0-constant},
\Cref{cube-thm-seesaw} and
\Cref{cube-cor-seesaw-trivial}.  Two of them are theorems and eight are
supporting lemmas; none of the eight is deep, and three of them ---
\Cref{cube-lem-collapse-neighbourhood},
\Cref{cube-lem-invertible-fibrewise} and
\Cref{cube-lem-projection-proper-flat} --- use nothing missing beyond the
definitions listed above, so they can be formalised today.

\subsection*{Summary}

Six definitions, of which four carry content and two are abbreviations over
what Mathlib already has; three propositions that cannot currently be stated
and are therefore left as open obligations, all three degree-zero consequences
of a single missing theory, the cohomology of sheaves of modules on schemes
and its behaviour in proper flat families; and ten statements proved in full
here, three of which are formalisable immediately.  Of the two theorems,
\Cref{cube-thm-rigidity} needs exactly one missing mathematical input, the
finite-dimensionality of $\Gamma(X,F)$ for $F$ coherent on a proper scheme
over a field, and \Cref{cube-thm-seesaw} needs that one together with
semicontinuity and cohomology and base change in degree zero.
